\documentclass[11pt,a4paper]{article}
\usepackage[T1]{fontenc}
\usepackage[utf8]{inputenc}
\usepackage{lmodern,amsmath,amssymb}
\usepackage[margin=25mm]{geometry}
\usepackage[hidelinks]{hyperref}
\hypersetup{pdftitle={The truncation error of a flow step is Gaussian times $e^{2t\|G\|_\infty}$, so the renorma},pdfauthor={navstokgap research working notes}}
\newcommand{\tightlist}{\setlength{\itemsep}{0pt}\setlength{\parskip}{0pt}}
\setlength{\emergencystretch}{3em}
\setcounter{secnumdepth}{0}
\begin{document}
\hypertarget{the-truncation-error-of-a-flow-step-is-gaussian-times-e2tg_infty-so-the-renormalization-step-is-small-exactly-on-small-field-configurations}{%
\section{\texorpdfstring{The truncation error of a flow step is Gaussian
times \(e^{2t\|G\|_\infty}\), so the renormalization step is small
exactly on small-field
configurations}{The truncation error of a flow step is Gaussian times e\^{}\{2t\textbackslash\textbar G\textbackslash\textbar\_\textbackslash infty\}, so the renormalization step is small exactly on small-field configurations}}\label{the-truncation-error-of-a-flow-step-is-gaussian-times-e2tg_infty-so-the-renormalization-step-is-small-exactly-on-small-field-configurations}}

The linearization of the gradient flow is the covariant heat equation
with a curvature term,
\[\partial_s\,\delta B_\mu=D^2\,\delta B_\mu+2\big[G_{\mu\nu},\delta B_\nu\big],
\qquad D_\mu=\partial_\mu+[B_\mu,\cdot\,],\] so the flow Jacobian
\(D\Phi_t\) is the kernel of that equation. Kato's inequality dominates
the covariant heat kernel by the free one pointwise, and the curvature
term is a bounded multiplication operator, so Duhamel gives
\[\big|D\Phi_t(x,y)\big|\ \le\ e^{2t\|G\|_\infty}\;\frac{1}{(4\pi t)^{3/2}}\,
e^{-|x-y|^2/(4t)} .\] Truncating the conjugated Hamiltonian of
\href{flow-conjugation-truncation.md}{the flow-conjugation note} at
range \(R\) therefore costs a relative error of order
\[\varepsilon(R,t)\ \sim\ \exp\Big[2t\|G\|_\infty-\frac{R^2}{4t}\Big],\]
and with \(R\) of order the scale \(\sqrt{8t}\) the exponent is
\(2t\|G\|_\infty-2\): the truncation is small \textbf{precisely when
\(t\,\|G\|_\infty\lesssim1\)}, that is when the flowed field strength is
below the inverse square of the scale being integrated out. This is the
small-field condition of constructive renormalization group, derived
here from the flow rather than assumed, and it explains why Balaban's
programme splits configurations into small-field and large-field
regions: on small fields the flow truncation is exponentially accurate,
and on large fields the Jacobian bound degrades exponentially and no
truncation of this kind is available. The flow's own monotonicity
controls \(\int|G|^2\) and leaves \(\|G\|_\infty\) open, so the missing
estimate is a pointwise bound on the flowed field strength. Constants
explicit; nothing promoted.

\begin{quote}
\textbf{Sharpness.} The factor \(e^{2t\|G\|_\infty}\) below cannot be
improved: the curvature operator is symmetric, its largest eigenvalue is
attained by the Nielsen--Olesen mode, and that mode grows under the flow
at exactly the rate \(2\|G\|\); see
\href{flow-instability-large-field.md}{the instability note}.
\end{quote}

\begin{quote}
\textbf{Corrected in part.} The statements below are continuum
statements at fixed physical scale. Within one lattice renormalization
step the compactness of the gauge group bounds \(a^2\|G\|_\infty\) by
\(\pi\) and the flow time of a doubling is \(a^2/2\), so the exponential
factor is a pure number and the truncation error is uniformly small; see
\href{lattice-truncation-uniform.md}{the lattice-truncation note}.
\end{quote}

\hypertarget{the-linearized-flow}{%
\subsection{1. The linearized flow}\label{the-linearized-flow}}

Take the continuum flow in the gauge-modified form
\(\partial_sB_\mu=D_\nu G_{\nu\mu}+\lambda D_\mu\partial_\nu B_\nu\)
with \(\lambda=1\) (Lüscher, arXiv:1006.4518v3, equations (1.1)--(1.2)
and (2.2); passage level via the
\href{../docs/Luscher_WilsonFlow_1006.4518v3.md}{local companion}).

\textbf{Proposition 1.} The variation \(\delta B\) of a solution obeys
\[\partial_s\,\delta B_\mu=D^2\,\delta B_\mu+2\big[G_{\mu\nu},\delta B_\nu\big].\]

\emph{Proof.} Varying the first term, \[\delta\big(D_\nu G_{\nu\mu}\big)
=D_\nu\big(D_\nu\delta B_\mu-D_\mu\delta B_\nu\big)+\big[\delta B_\nu,G_{\nu\mu}\big]
=D^2\delta B_\mu-D_\nu D_\mu\delta B_\nu+\big[\delta B_\nu,G_{\nu\mu}\big],\]
and
\(D_\nu D_\mu\delta B_\nu=D_\mu D_\nu\delta B_\nu+[G_{\nu\mu},\delta B_\nu]\).
Varying the gauge term gives \(D_\mu D_\nu\delta B_\nu\) to the same
order, which cancels the second piece, leaving
\(D^2\delta B_\mu-2[G_{\nu\mu},\delta B_\nu]=D^2\delta B_\mu+2[G_{\mu\nu},\delta B_\nu]\).
\(\square\)

So \(D\Phi_t\) is the solution kernel of a covariant heat equation with
a zeroth-order term given by the curvature acting in the adjoint
representation.

\hypertarget{the-gaussian-bound}{%
\subsection{2. The Gaussian bound}\label{the-gaussian-bound}}

\textbf{Proposition 2.} Let
\(\|G\|_\infty=\sup_{0\le s\le t,\ x}\|G(s,x)\|\) in the adjoint
operator norm. Then the kernel of the linearized flow obeys
\[\big|D\Phi_t(x,y)\big|\ \le\ e^{2t\|G\|_\infty}\,K^{\rm free}_t(x-y),
\qquad K^{\rm free}_t(z)=\frac{e^{-|z|^2/(4t)}}{(4\pi t)^{3/2}} .\]

\emph{Proof.} Write the equation as \(\partial_s u=D^2u+Mu\) with
\((Mu)_\mu=2[G_{\mu\nu},u_\nu]\), a multiplication operator of norm at
most \(2\|G\|_\infty\). Duhamel gives
\(u(t)=e^{tD^2}u(0)+\int_0^t e^{(t-s)D^2}Mu(s)\,ds\), and iterating,
\(\|u(t)\|\le\sum_n\frac{(2t\|G\|_\infty)^n}{n!}\sup\|e^{\cdot D^2}\|\)-type
bounds; pointwise, Kato's inequality gives
\(|e^{sD^2}f|\le e^{s\Delta}|f|\) for the covariant Laplacian in any
representation (Simon, J. Funct. Anal. 32 (1979) 97, metadata level), so
each term is dominated by the free heat kernel and the series sums to
\(e^{2t\|G\|_\infty}\). \(\square\)

The bound is saturated in the free case, where \(G=0\) and
\(D\Phi_t=e^{t\Delta}\) exactly, in agreement with
\href{flowed-bound-free-field.md}{the free-field note} §2.

\hypertarget{the-truncation-error}{%
\subsection{3. The truncation error}\label{the-truncation-error}}

The conjugated kinetic operator of
\href{flow-conjugation-truncation.md}{the flow-conjugation note} §2 is
built from \(D\Phi_t\), which couples links at separation \(|x-y|\) with
weight bounded by Proposition 2. Truncating it to range \(R\) discards
the tail \(|x-y|>R\), whose weight relative to the whole is
\[\varepsilon(R,t)\ \le\ e^{2t\|G\|_\infty}\int_{|z|>R}K^{\rm free}_t(z)\,d^3z
\ \le\ C\,e^{2t\|G\|_\infty}\,e^{-R^2/(4t)}\Big(1+\frac{R}{\sqrt t}\Big),\]
by the standard Gaussian tail estimate. With the natural choice
\(R=\sqrt{8t}\), the exponent is \(2t\|G\|_\infty-2\), so

\textbf{Corollary 3.} The flow-truncation step at scale \(\sqrt{8t}\)
has relative error
\[\varepsilon\ \lesssim\ \exp\big[2t\|G\|_\infty-2\big],\] small when
\(t\|G\|_\infty\lesssim1\) and exponentially large otherwise. Taking
\(R=\kappa\sqrt{8t}\) improves the exponent to
\(2t\|G\|_\infty-2\kappa^2\) at the cost of a longer-range effective
Hamiltonian, so a fixed accuracy is available whenever
\(\kappa^2\gtrsim t\|G\|_\infty\).

\hypertarget{what-the-condition-says}{%
\subsection{4. What the condition says}\label{what-the-condition-says}}

Write \(\ell=\sqrt{8t}\) for the scale being integrated out. The
condition \(t\|G\|_\infty\lesssim1\) reads
\[\|G\|_\infty\ \lesssim\ \frac{8}{\ell^2},\] that is, the field
strength must be below the inverse square of the scale. In the units of
\href{mass-gap-obligations-lattice.md}{the obligations map} the lattice
field strength at spacing \(a\) is at most of order \(1/a^2\), so at the
first step, \(\ell\sim a\), the condition is marginal; after the field
has been smoothed over \(\ell\gg a\) the typical \(\|G\|\) falls and the
condition is comfortable on smooth configurations and violated on rough
ones.

This is the small-field condition of constructive renormalization group,
obtained here as the condition for the flow truncation to be accurate,
and it identifies the difficulty by name:

\begin{itemize}
\tightlist
\item
  on \textbf{small-field} configurations, \(\|G\|_\infty\le\eta/\ell^2\)
  with \(\eta\) small, the truncation error is \(e^{-2\kappa^2+2\eta}\)
  and every step of the route of
  \href{flow-conjugation-truncation.md}{the flow-conjugation note} §6 is
  controlled;
\item
  on \textbf{large-field} configurations the Jacobian bound degrades
  like \(e^{2t\|G\|_\infty}\) and no truncation of this kind is
  available, so those regions must be handled by a different argument.
\end{itemize}

The split is exactly the one Balaban's programme makes (Commun. Math.
Phys. 122 (1989) 355 and the accompanying series, metadata level), and
the derivation above says why it is forced rather than chosen.

\hypertarget{what-the-flow-itself-controls-and-what-it-leaves}{%
\subsection{5. What the flow itself controls, and what it
leaves}\label{what-the-flow-itself-controls-and-what-it-leaves}}

The flow decreases the action monotonically,
\(\frac{d}{ds}S(B_s)=-\|D^*G\|_2^2\le0\)
(\href{comparison-and-bridges.md}{comparison note} §4), so
\[\int|G(t,x)|^2\,d^3x\ \le\ \int|G(0,x)|^2\,d^3x\] for all \(t\): an
\(L^2\) bound, uniform in the flow time. Proposition 2 needs an
\(L^\infty\) bound, and the gap between them is the whole question. Two
remarks fix its size.

\emph{The free case is a parabolic smoothing estimate.} For \(G=0\) the
linearized flow is the heat semigroup and
\(\|e^{t\Delta}f\|_\infty\le(4\pi t)^{-3/4}\|f\|_2\), so an \(L^2\)
bound on the initial curvature gives
\(\|G(t)\|_\infty\lesssim t^{-3/4}\|G(0)\|_2\), which beats the required
\(1/t\) for large \(t\) and fails for small \(t\). The interacting
statement of this type is the a priori estimate that a construction must
supply.

\emph{Monotonicity is not enough by itself.} An \(L^2\) bound permits a
configuration with curvature concentrated on a small region at arbitrary
height, which is precisely a large-field region; excluding it is a
statement about the measure in the Euclidean formulation and about the
state in the Hamiltonian one. In the Hamiltonian framework there is no
integration over configurations to absorb such regions into a small
probability, which makes the Hamiltonian route harder at exactly this
point than the Euclidean one.

\hypertarget{consequence-for-state}{%
\subsection{6. Consequence for STATE}\label{consequence-for-state}}

The truncation error of a flow-based renormalization step is now
explicit: Gaussian in the ratio of the truncation range to the flow
radius, multiplied by \(e^{2t\|G\|_\infty}\). The step is controlled
exactly on configurations with \(\|G\|_\infty\lesssim\ell^{-2}\), which
derives the small-field condition instead of assuming it. The missing
estimate is pointwise control of the flowed field strength, an
\(L^\infty\) bound where the flow supplies only \(L^2\), and in the
Hamiltonian framework there is no probabilistic route around it. That is
the sharpest statement of the remaining obstacle this programme has
reached, and it points at the Euclidean formulation as the place where
the large-field regions can be given small weight.
\end{document}
